\section*{1.16\quad Энергия}
Мы приступаем к изучению энер\-гии — фундаментального физического понятия. Но предвари\-тельно нужно разобраться с другой физической величиной — работой силы.

\kern-.149bp\relax
\subsection*{1.16.1\quad Работа}
Пусть на тело действует постоянная сила $\vec F$ и тело, двигаясь прямолинейно по горизонтальной поверхности, совершило перемещение $\vec s$. Сила $\vec F$ не обязательно является непосредственной при\-чиной перемещения (так, сила тяже\-сти не является непосредственной при\-чиной перемещения шкафа, который передвигают по комнате).

Предположим сначала, что векторы силы и перемещения сонаправлены (рис.~\textcolor{booklink}{1.48}; остальные силы, действующие на тело, не указаны).
\ReplicaWorkFigureOne

В этом простейшем случае работа $A$ определяется как произведение модуля силы на модуль перемещения:
\begin{equation}A=Fs.\tag{1.76}\end{equation}
\vskip-.432bp\relax


Единицей измерения работы служит \textit{джоуль} (Дж): Дж = Н $\cdot$ м. Таким образом, если под действием силы 1 Н тело перемещается на 1 м, то сила совершает работу 1 Дж.

Работа силы, перпендикулярной перемещению, по определению считается равной нулю. Так, в данном случае сила тяже\-сти и сила реакции опоры не соверша\-ют работы.

Пусть теперь вектор силы образует с вектором перемещения острый угол $\alpha$ (рис.~\textcolor{booklink}{1.49}).
\ReplicaWorkFigureTwo

Разложим силу $\vec F$ на две составляющие: $\vec F_{\parallel}$ (параллельную перемещению) и $\vec F_{\perp}$ (перпенди\-кулярную перемещению). Работу совершает только $\vec F_{\parallel}$. Поэтому для работы силы $\vec F$ получаем: $A=F_{\parallel}s=F\cos\alpha\cdot s$. Итак,
\begin{equation}A=Fs\cos\alpha.\tag{1.77}\end{equation}
\vskip-.432bp\relax


Если вектор силы образует с вектором перемещения тупой угол $\alpha$, то работа по-прежнему определяется формулой (\textcolor{booklink}{1.77}). В этом случае работа оказывается отрицательной.

Например, работа силы трения скольжения, действующей на тело в рассмотренных ситуа\-циях, будет отрицательной, так как сила трения направлена противоположно перемещению. В этом случае имеем: $\alpha=180^{\circ}$, $\cos\alpha=-1$, и для работы силы трения получаем:
\vadjust{\kern-1.007bp}
$$A_{\text{тр}}=-F_{\text{тр}}s=-\mu mgs,$$
\newpage
\noindent где $m$ — масса тела, $\mu$ — коэффициент трения между телом и опорой.

\vspace{5.146bp}
Соотношение (\textcolor{booklink}{1.77}) означает, что работа является \textit{скалярным произведением} векторов силы и перемещения:
$$A=\vec F\vec s.$$

Это позволяет вычислять работу через координаты данных векторов:
$$A=F_xs_x+F_ys_y+F_zs_z.$$

Пусть на тело действуют несколько сил $\vec F_1,\vec F_2,\ldots,\vec F_n$, и $\vec F$ — равнодействующая этих сил. Для работы силы $\vec F$ имеем:
$$A=\vec F\vec s=(\vec F_1+\vec F_2+\ldots+\vec F_n)\vec s=\vec F_1\vec s+\vec F_2\vec s+\ldots+\vec F_n\vec s,$$
\noindent или
$$A=A_1+A_2+\ldots+A_n,$$
\noindent где $A_1,A_2,\ldots,A_n$ — работы сил $\vec F_1,\vec F_2,\ldots,\vec F_n$. Итак, \textit{работа равнодействующей приложенных к телу сил равна сумме работ каждой силы в отдельности.}

\subsection*{1.16.2\quad Мощность}
Часто имеет значение быстрота, с ко\-торой совершается работа. Скажем, на практике важно знать, какую работу сможет выполнить данное устройство за фиксированное время.

\textit{Мощность} — это величина, характеризующая скорость совершения работы. Мощность $N$ есть отношение работы $A$ ко времени $t$, за которое эта работа совершена:
$$N=\frac{A}{t}.$$

Мощность измеряется в \textit{ваттах} (Вт). 1 Вт = 1 Дж/с, то есть 1 Вт — это такая мощность, при ко\-торой работа в 1 Дж совершается за 1 с.

Предположим, что силы, действующие на тело, уравновешены, и тело движется равномерно и прямолинейно со скоростью $\vec v$. В этом случае существует полезная формула для мощности, развиваемой одной из действующих сил $\vec F$.

За время $t$ тело совершит перемещение $\vec s=\vec vt$. Работа силы $\vec F$ будет равна:
$$A=\vec F\vec s=\vec F\vec vt.$$
\noindent Отсюда получаем мощность:
$$N=\vec F\vec v,$$
\noindent или
$$N=Fv\cos\alpha,$$
\noindent где $\alpha$ — угол между векторами силы и скорости.

Наиболее часто эта формула используется в ситуации, когда $\vec F$ — «сила тяги» двигателя автомобиля (которая на самом деле есть сила трения ведущих колёс о дорогу). В этом случае $\alpha=0$, и мы получаем просто:
$$N=Fv.$$
\newpage
\subsection*{1.16.3\quad Механическая энергия}
Энергия является мерой движе\-ния и взаимодействия любых объектов в природе. Имеются различные формы энер\-гии: механиче\-ская, тепловая, электромагнитная, ядерная\ldots

Опыт показывает, что энергия не появляется ниоткуда и не исчезает бесследно, она лишь переходит из одной формы в другую. Это самая общая формулировка \textit{закона со\-хране\-ния энер\-гии.}

Каждый вид энер\-гии представляет собой некоторое математическое выражение. Закон со\-хране\-ния энер\-гии означает, что в каждом явлении природы определённая сумма таких выра\-жений остаётся постоянной с течением времени.

Измеряется энергия в джоулях, как и работа.

\textit{Механическая энергия} является мерой движе\-ния и взаимодействия механических объектов (материальных точек, твёрдых тел).

Мерой движе\-ния тела является \textit{кинетическая энергия}. Она зависит от скорости тела. Мерой взаимодействия тел является \textit{потенциаль\-ная энергия}. Она зависит от взаимного расположения тел.

Механическая энергия си\-стемы тел равна сумме кинетической энер\-гии тел и потенциаль\-ной энер\-гии их взаимодействия друг с другом.

\kern-.487bp\relax
\subsection*{1.16.4\quad Кинетическая энергия}
Кинетической энергией тела (принимаемого за материальную точку) называется величина
\vskip-2.801bp\relax
$$K=\frac{mv^2}{2}\,,$$
\vskip2.181bp\relax
\noindent где $m$ — масса тела, $v$ — его скорость.

Кинетической энергией си\-стемы из $N$ тел называется сумма кинетических энергий каждого тела:
$$K=\frac{m_1v_1^2}{2}+\frac{m_2v_2^2}{2}+\ldots+\frac{m_Nv_N^2}{2}\,.$$
\vskip-1.413bp\relax

Если тело движется под действием силы $\vec F$, то кинетическая энергия тела, вообще гово\-ря, ме\-няется со временем. Оказывается, именение кинетической энер\-гии тела за некоторый промежуток времени равно работе силы $\vec F$. Покажем это для случая прямолинейного равно\-ускоренного движе\-ния.

Пусть $\vec v_1$ — начальная скорость, $\vec v_2$ — конечная скорость тела. Выберем ось $X$ вдоль тра\-ек\-тории тела (и, соответственно, вдоль вектора силы $\vec F$). Для работы силы $\vec F$ получаем:
\vskip-2.801bp\relax
$$A=\vec F\vec s=F_xs_x=ma_xs_x=ma_x\frac{v_{2x}^2-v_{1x}^2}{2a_x}=\frac{mv_{2x}^2-mv_{1x}^2}{2}\,.$$
\vskip2.181bp\relax
\noindent (мы воспользовались формулой для $s_x$, выведенной в разделе «Равноускоренное движение»).

Заметим теперь, что в данном случае проекция скорости отличается от модуля скорости разве что знаком; поэтому $v_{1x}^2=v_1^2$ и $v_{2x}^2=v_2^2$. В результате имеем:
\vskip-2.833bp\relax
$$A=\frac{mv_2^2}{2}-\frac{mv_1^2}{2}=K_2-K_1=\Delta K,$$
\vskip2.181bp\relax
\noindent что и требовалось.

На самом деле соотношение $\Delta K=A$ справедливо и в самом общем случае криволинейного движе\-ния под действием переменной силы.

\vskip5.146bp\relax
\noindent\textbf{Теорема о кинетической энер\-гии.} Изменение кинетической энер\-гии тела равно работе, со\-вершённой приложенными к телу внешними силами за рассматриваемый промежуток времени.
\newpage
